\answer{ex:22:1}{(a)~$32\%$ と $100\%$。(b)~$R_\uparrow-R_\downarrow\ge2\sqrt{2000\,\text{Hz}/1\,\text{s}}\approx89$~Hz。合計頻度のわずか $4.5\%$ である。}
\answer{ex:22:2}{(a)~$n\ge4(p_1q_1+p_2q_2)/\Delta^2\approx560$。(b)~$\approx56$。}
\answer{ex:22:3}{(a)~詳細釣り合い：$\rho PK$ は対称である。(b)~ミスの割合は調和平均 $1/\langle1/P\rangle=0.18$ で、フィードバックなしの $0.5$ に比べて小さい。}
\answer{ex:22:4}{(a)~面積 $4h^2=0.04$ の平行四辺形（$p$ の制限を考慮すると数値的に $0.045$）。(b)~$\approx0.18$。}
\answer{ex:22:5}{$0.49$：一部の培養は、ほぼ必ずミスする領域に入り込み、そこをさまよう。制限をつけると $1.00$：有界な集合では、密度 $\propto1/P$ は吸収領域に集まる。}
\answer{ex:22:6}{(a)~$\ge5$ 段階が必要で、$8$ 本の電極で足りる。(b)~不可能：$r_{\max}\ge25/T=100$~Hz が必要になる。}
\answer{ex:22:7}{未解決の課題。ランダム探索の時間は $d$ とともにおよそ体積比 $(L/h)^d$ のように増え、誤差の符号を使う探索では $d$ に比例して増える。仮説を区別するのは入れ替え実験である。ミスの後に予測可能だが強い刺激、ヒットの後に予測不能だが弱い刺激を与える。各群に数十個の培養が必要である。}
